\section{Completion of the reduction}
\label{sec:finish}

The matrix game has a fixed soundness gap and an arbitrarily small
completeness error.  We now obtain the form required by
Theorem~\ref{thm:ugc}: unweighted, simple bipartite games with arbitrary
prescribed errors.  Rounding the occurrence weights and subdividing
each occurrence will preserve a fixed gap.  One application of
alphabet-independent repetition then completes the proof.

\subsection{Removing weights and repeated edges}

We use the following elementary construction for weighted permutation
constraints.  At this intermediate stage, a constraint occurrence may
be a loop, and several occurrences may join the same vertices.

\begin{lemma}[Rounding and subdivision]
\label{lem:unweighted-subdivision}
Let $H$ be a finite permutation-constraint game with $h\ge1$ positive
rational occurrence weights summing to one.  For every positive
rational $\tau$, put $Q=\lceil h/\tau\rceil$.  There is a
deterministic construction of an unweighted simple bipartite game $B$
with $4Q$ edges and the same alphabet such that
\begin{equation}
 \left|\val(B)-\left(1-\frac{1-\val(H)}4\right)\right|
 \le\frac{\tau}{4}.
 \label{eq:rounded-subdivision-value}
\end{equation}
If the constraints of $H$ are translations on a finite binary vector
space, so are those of $B$.  For fixed $\tau$, when $H$ is listed by its
vertices, occurrences, permutation tables, and binary rational weights,
the construction takes polynomial time in its encoding length.
\end{lemma}

\begin{proof}
Write the occurrence weights as $w_1,\ldots,w_h$.  Start with
$n_e=\lfloor Qw_e\rfloor$, and add one to the
$Q-\sum_e\lfloor Qw_e\rfloor$ entries having largest fractional parts,
breaking ties in a fixed order.  Then
\[
 n_e\ge0,\qquad \sum_e n_e=Q,\qquad
 |w_e-n_e/Q|\le1/Q.
\]
The weight differences sum to zero.  Consequently, for every labeling,
the change in the weight of its satisfied-occurrence set is at most
\[
 \frac12\sum_{e=1}^h|w_e-n_e/Q|
 \le\frac{h}{2Q}\le\tau.
\]
Replace each occurrence $e$ by $n_e$ copies, omitting it if $n_e=0$.
The resulting unweighted multigraph $H'$ has $Q$ occurrences and
$|\val(H')-\val(H)|\le\tau$.

Orient each occurrence $e$ of $H'$ from $u$ to $v$, writing its
constraint as $a_v=\rho_e(a_u)$.  Replace it by the path
\[
 u\;--\;p_e\;--\;q_e\;--\;r_e\;--\;v,
\]
using three fresh interior vertices for \emph{every occurrence}.
Place all original vertices and all $q_e$ on the left, and all
$p_e,r_e$ on the right.  In the direction from $u$ to $v$, put identity
permutations on the first three edges and $\rho_e$ on the last.
When expressing an edge in left-to-right orientation, invert its
permutation if needed.  Give all $4Q$ edges equal weight.

Fix labels at the original vertices.  A satisfied occurrence extends
to four satisfied edges.  A violated occurrence cannot extend to four,
but propagating labels along the first three edges satisfies exactly
three.  Since the interiors are disjoint, these extensions can be
chosen independently.  Maximizing over the original labels gives the
exact identity
\begin{equation}
 \val(B)=1-\frac{1-\val(H')}{4}.
 \label{eq:bipartization-value}
\end{equation}
This proves Equation~\eqref{eq:rounded-subdivision-value}.

A loop becomes a four-cycle with three distinct new vertices.
Parallel occurrences have disjoint interiors.  Thus each left--right
vertex pair carries at most one edge, so $B$ is simple.  An inverse
binary translation is the same translation, and an identity is
translation by zero; the claimed constraint form is preserved.

Finally, $Q\le h/\tau+1$ is polynomially bounded for fixed $\tau$.
Integer division and exact comparison of rational remainders compute
the $n_e$ in polynomial time and bit complexity.  Listing their $Q$
copies and the $3Q$ new vertices is therefore polynomial as well.
\end{proof}

\subsection{The final reduction}

\begin{proof}[Proof of Theorem~\ref{thm:ugc}]
Fix $\eps,\delta\in(0,1/2)$.  Choose positive rational
$\eps_0\le\eps$ and $\delta_0\le\delta$.  We will take
$\tau\le1/200$ when applying the rounding and subdivision lemma.
For any matrix game $H$ satisfying the NO bound
$\val(H)\le99/100$, the resulting game $B$ then satisfies
\begin{equation}
 \val(B)\le1-\frac{1/100-\tau}{4}\le1-\frac1{800}.
 \label{eq:post-subdivision-gap}
\end{equation}
This gap does not depend on the alphabet.  We can therefore use
alphabet-independent repetition to choose its exponent before the
latent construction fixes that alphabet.  Choose an integer $t\ge1$
such that
\begin{equation}
 \left(1-\frac1{10240000}\right)^t\le\delta_0.
 \label{eq:final-repetition-gap}
\end{equation}
For example, $t\ge\lceil10240000\log(1/\delta_0)\rceil$ suffices.
By Theorem~\ref{thm:repetition}, this exponent reduces every projection
game of value at most $1-1/800$ to value at most $\delta_0$.

Choose positive rational parameters
\begin{equation}
 p_*\le\frac{\eps_0}{t},\qquad
 \tau\le\min\left\{\frac1{200},\frac{\eps_0}{t}\right\}.
 \label{eq:final-errors}
\end{equation}
Choose the latent rank threshold $r_*$ from Lemma~\ref{lem:latent},
put $\eta=2^{-r_*}/16$, and choose the corresponding shortcode
parameters $\alpha,r_{\rm s}$, using a positive rational lower bound
for $\alpha$.  Then choose a sufficiently large $\ell$, as specified
in Equation~\eqref{eq:matrix-parameters}, and fix the corresponding
latent data $(\mathcal V,K,C,\mu)$ and a splitting
$\mathcal V=K\oplus W$, with $K=\F_2^\ell$.

Once the rational budgets are fixed, these choices are effective.
For rational $p_*$, the effective assertion of Lemma~\ref{lem:latent}
supplies all the finite latent data, including a computable rational
noise law, by a terminating deterministic construction.  The effective bounds stated after
Theorem~\ref{thm:shortcode} supply the rational choice of $\alpha$ and
the dimension thresholds, including $m_0(\eta,\ell)$ after $\ell$ has
been fixed.

Next choose a cube $k\ge8$ large enough for
Proposition~\ref{prop:matrix-gap}.
Such a cube can be found by checking the estimates used in its proof.
Write $m=1+2k$ and $\beta=k^{-2/3}$; since $k$ is a cube, $\beta$ is
rational.  Require $m\ge m_0(\eta,\ell)$ and
\[
 \frac{\alpha^2}{32}\,2^{(\ell-r_{\rm s})(m-r_{\rm s})}
 >(2r_{\rm s}+1)(\ell+m).
\]
The latter is a sufficient rational condition for the union bound in
Lemma~\ref{lem:good-row-advice} to be less than one.  For witness
transfer in Lemma~\ref{lem:projected-witness}, require
\[
 \frac14\left[
 \left(1+\frac{|\mathcal V|-1}{3}\beta^2\right)^k-1\right]
 <\frac{\gamma_1^2}{(1+\alpha/4)^2},
\]
where $\gamma_1$ is from
Equation~\eqref{eq:projected-witness-mass}.  The left-hand side is the
square of the total-variation bound.  For the final contradiction,
require
\[
 \left(1-\frac{\beta\,2^{-\dim W-r_{\rm s}}}{3600}\right)^k
 <\gamma,
\]
with $\gamma$ from Proposition~\ref{prop:matrix-decoding}.  All these
conditions are rational and hold for sufficiently large cubes $k$.
Enumerating cubes until they hold therefore finds a suitable $k$;
the search cost depends only on the two prescribed errors and is
incurred before the input formula is used.

Finally, with $k$ fixed, choose a positive rational
\begin{equation}
 \xi<\min\left\{\frac1{100},\frac{\eps_0}{2kt}\right\}.
 \label{eq:final-parity-error}
\end{equation}
Every choice depends only on the two prescribed errors.  In
particular, $r_*$ is fixed before $\ell$, the advice dimensions before
$k$, and $k$ before the parity error $\xi$.

Given a 3SAT formula, apply Theorem~\ref{thm:parity} and
Lemma~\ref{lem:distinct-parity}, and construct the matrix game $H$
of Section~\ref{sec:game}.  Proposition~\ref{prop:matrix-gap} gives
\[
 \begin{array}{ll}
 \text{YES:}&\val(H)\ge1-d,
       \quad d=k\xi+p_*/2\le\eps_0/t,\\[2pt]
 \text{NO:}&\val(H)\le99/100.
 \end{array}
\]
Apply Lemma~\ref{lem:unweighted-subdivision} with $\tau$ to obtain
the unweighted simple bipartite game $B$.  It satisfies
\[
 \begin{array}{ll}
 \text{YES:}&\val(B)\ge1-(d+\tau)/4,\\[2pt]
 \text{NO:}&\val(B)\le1-1/800,
 \end{array}
\]
where the NO bound is Equation~\eqref{eq:post-subdivision-gap}.
The final output is $G=B^{\otimes t}$.  Its NO value is at most
$\delta_0\le\delta$, by Equations~\eqref{eq:projection-repetition}
and \eqref{eq:final-repetition-gap}.  On a YES input, use a labeling
of $B$ coordinatewise.  A union bound gives
\[
 \val(G)\ge1-\frac{t(d+\tau)}4
 \ge1-\eps_0/2\ge1-\eps.
\]

The output has the claimed graph and constraint form.  Its two vertex
sets are the $t$-fold products of the two sides of $B$.  A pair of
product vertices determines at most one edge in each coordinate,
since $B$ is simple, and hence at most one product edge.  Independent
uniform edge sampling is uniform on these product edges.  The
constraints of $H$ are translations on $K$, by
Section~\ref{sec:game}; rounding and subdivision preserve them.
The product of translations with offsets $c_1,\ldots,c_t$ is
translation by $(c_1,\ldots,c_t)$ on $K^t=\F_2^{\ell t}$.  This is
the alphabet $\F_2^s$ in the theorem, with $s=\ell t\ge1$ fixed by
the prescribed errors.

\paragraph{Polynomial encoding.}
Suppose the preprocessed parity instance has $N_{\mathcal E}$ positive
equation occurrences.  There are $N_{\mathcal E}^k$ equation tuples.
The number of tables, noise outcomes, and perturbation vectors for each
tuple is fixed after the parameter choices above, so the number $h$ of
listed matrix-test occurrences is a fixed constant times
$N_{\mathcal E}^k$.  Omit
zero-probability outcomes and retain repeated occurrences separately.
Each weight is a product of $k$ parity weights, a fixed rational
noise probability, and fixed uniform-sampling factors.  It has
polynomial binary length.

Exact table keys and folded representatives are computed by linear
algebra in fixed dimensions, keeping the original variable and
occurrence IDs.  At most $2h$ endpoints appear.  Even comparing every
pair of keys to identify equal vertices takes polynomial time; number
the resulting distinct vertices consecutively.
Largest-remainder rounding uses $Q=\lceil h/\tau\rceil=O(h)$ copies;
subdivision creates $3Q$ new vertices and $4Q$ edges.  The final
product has $(4Q)^t$ edges and polynomially many vertices, because
$t$ is fixed.  Full permutation tables on the fixed alphabet
$\F_2^{\ell t}$ have constant size, and endpoint indices have
$O(t\log(h+1))$ bits.  These observations give a deterministic
polynomial-time explicit construction throughout and complete the
proof.
\end{proof}
